
\section{BETTER PRACTICES FOR BODY BIASING INJECTION}
%\begin{frame}
%    \frametitle{Why working with BBI?}
%    \begin{itemize}
%        \setlength\itemsep{1em}
%        \item BBI introduced to perform Bellcore attack;
%        \item Demanding fault attacks difficult or impossible to perform;
%        \item Low experiment repeatability.
%    \end{itemize}
%    \vspace{1cm}
%    \begin{center}
%        \huge What are the limiting factors?
%    \end{center}
%\end{frame}

\begin{frame}
    \frametitle{State-of-the-art BBI platform: limiting factors}
    \begin{columns}
        \begin{column}{0.65\textwidth}
            \centering
            \begin{figure}
                \includegraphics[width=0.6\textwidth]{model0Probe.pdf}
            \end{figure}
%            \vspace{-0.5cm}
            \begin{figure}
                \includegraphics[width=0.45\textwidth]{S1P.pdf}
                \includegraphics[width=0.45\textwidth]{S1G.pdf}
            \end{figure}
        \end{column}
        \barsep
        \begin{column}{0.35\textwidth}
            \centering
            \begin{itemize}
                \setlength\itemsep{1em}
                \item Impedance mismatch;
                \item Floating grounds;
                \item Ringing;
                \item Set-point error.
            \end{itemize}
        \end{column}
    \end{columns}
\end{frame}
